% from ega4-21a-fr.tex
\oldpage[IV-4]{259}
multiplication par ces éléments fournit tous les automorphismes du $\sh{O}_U$-Module $\sh{O}_U$ (\hyperref[0.5.4.3]{$\mathbf{0}_{\mathrm I}$, 5.4.3}).

\begin{corollary}[21.2.3]
\label{IV.21.2.3}
\begin{enumerate}[label=(\roman*)]
  \item Soit $\sh{J}$ un Idéal fractionnaire inversible;
  alors le $\sh{O}_X$-Module inversible $\sh{J}^{-1}$ s'identifie canoniquement à l'Idéal fractionnaire $\sh{J}'$ (\emph{transporteur} de $\sh{J}$ dans $\sh{O}_X$) défini de la façon suivante : pour tout ouvert $U$ de $X$ tel que $\sh{J}|U=\sh{O}_U.f$, où $f\in\Gamma(U,\sh{M}_X^*)$, on a $\sh{J}'|U=\sh{O}_U.f^{-1}$.
  \item Si $\sh{J}_1$ et $\sh{J}_2$ sont deux Idéaux fractionnaires inversibles, l'application canonique $\sh{J}_1\otimes\sh{J}_2\to\sh{J}_1\sh{J}_2$ est un isomorphisme de $\sh{O}_X$-Modules.
\end{enumerate}
\end{corollary}

\begin{proof}
L'assertion (ii) résulte aussitôt de \sref{IV.21.2.2}.
D'autre part, la remarque faite à la fin de \sref{IV.21.2.2} prouve qu'il existe bien un Idéal fractionnaire $\sh{J}'$ et un seul défini par la condition de l'énoncé;
l'isomorphisme canonique de $\sh{J}'$ sur $\sh{J}^{-1}=\shHom_{\sh{O}_X}(\sh{J},\sh{O}_X)$ s'obtient en faisant correspondre à toute section $s(f^{-1}|V)$ de $\Gamma(V,\sh{J}')$ (où $V$ est un ouvert contenu dans $U$ et $s\in\Gamma(V,\sh{O}_X)$) l'homomorphisme $t(f|V)\mapsto st$ de $\Gamma(V,\sh{J})$ dans $\Gamma(V,\sh{O}_X)$.
\end{proof}

En vertu de \sref{IV.21.2.3},~(i), on identifiera généralement les $\sh{O}_X$-Modules inversibles $\sh{J}'$ et $\sh{J}^{-1}$, considérant donc $\sh{J}^{-1}$ comme un sous-$\sh{O}_X$-Module de $\sh{M}_X$.

\begin{env}[21.2.4]
\label{IV.21.2.4}
Il résulte de \sref{IV.21.2.3} que l'ensemble $\operatorname{Id.inv}(X)$ des Idéaux fractionnaires inversibles sur $X$ est muni d'une structure de \emph{groupe commutatif} pour la loi de composition $(\sh{J}_1,\sh{J}_2)\mapsto\sh{J}_1\sh{J}_2$, l'élément neutre de ce groupe étant $\sh{O}_X$.
Il est clair que pour tout ouvert $U$ de $X$, on a $(\sh{J}_1\sh{J}_2)|U=(\sh{J}_1|U)(\sh{J}_2|U)$, donc l'application de restriction $\sh{J}\mapsto\sh{J}|U$ est un homomorphisme de groupes $\operatorname{Id.inv}(X)\to\operatorname{Id.inv}(U)$;
on définit donc ainsi un \emph{préfaisceau de groupes commutatifs} $U\mapsto\operatorname{Id.inv}(U)$;
il est immédiat qu'en fait ce préfaisceau est un \emph{faisceau de groupes commutatifs}, que l'on note $\sh{J}d.inv_X$.
\end{env}

\begin{env}[21.2.5]
\label{IV.21.2.5}
Pour toute fonction méromorphe régulière $f\in\Gamma(X,\sh{M}_X^*)$, il résulte de \sref{IV.21.2.2} que $\sh{J}(f)=\sh{O}_X.f$ est un Idéal fractionnaire inversible, et on a évidemment $\sh{J}(f_1f_2)=\sh{J}(f_1)\sh{J}(f_2)$, autrement dit l'application $f\mapsto\sh{J}(f)$ est un homomorphisme du groupe commutatif $\Gamma(X,\sh{M}_X^*)$ dans le groupe commutatif $\operatorname{Id.inv}(X)$.
Remplaçant $X$ par un ouvert quelconque $U$ de $X$ et notant que les homomorphismes obtenus sont compatibles avec les opérations de restriction, on obtient un homomorphisme canonique de faisceaux de groupes commutatifs :
\[
\label{IV.21.2.5.1}
  I_0:\sh{M}_X^*\to\sh{J}d.inv_X.
  \tag{21.2.5.1}
\]
Notons que si $f\in\Gamma(X,\sh{O}_X^*)$, on a $\sh{J}(f)=\sh{O}_X$;
on en déduit aussitôt que l'homomorphisme $I_0$ se factorise en
\[
\label{IV.21.2.5.2}
  \sh{M}_X^*\to\sh{M}_X^*/\sh{O}_X^*=\sh{D}iv_X\xrightarrow{I}\sh{J}d.inv_X
  \tag{21.2.5.2}
\]
où $I$ est un homomorphisme du faisceau de groupes additifs $\sh{D}iv_X$ dans le faisceau de groupes multiplicatifs $\sh{J}d.inv_X$;
on a donc pour tout ouvert $U$ de $X$ un homomorphisme $\sh{J}_U:\operatorname{Div}(U)\to\operatorname{Id.inv}(U)$ de groupes commutatifs, tel que, pour toute section $f\in\Gamma(U,\sh{M}_X^*)$, on ait
\[
\label{IV.21.2.5.3}
  \sh{J}_U(\operatorname{div}_U(f))=\sh{O}_U.f.
  \tag{21.2.5.3}
\]

